% Options for packages loaded elsewhere
\PassOptionsToPackage{unicode}{hyperref}
\PassOptionsToPackage{hyphens}{url}
\PassOptionsToPackage{dvipsnames,svgnames,x11names}{xcolor}
\documentclass[
  polish,
  12pt,
  a4paper,
]{article}
\usepackage{xcolor}
\usepackage[margin=2.3cm]{geometry}
\usepackage{amsmath,amssymb}
\usepackage{fancyhdr}
\usepackage{titlesec}
\setcounter{secnumdepth}{-\maxdimen} % remove section numbering
\usepackage{iftex}
\ifPDFTeX
  \usepackage[T1]{fontenc}
  \usepackage[utf8]{inputenc}
  \usepackage{textcomp} % provide euro and other symbols
\else % if luatex or xetex
  \usepackage{unicode-math} % this also loads fontspec
  \defaultfontfeatures{Scale=MatchLowercase}
  \defaultfontfeatures[\rmfamily]{Ligatures=TeX,Scale=1}
\fi
\usepackage{lmodern}
\ifPDFTeX\else
  % xetex/luatex font selection
  \setmainfont[]{Times New Roman}
\fi
% Use upquote if available, for straight quotes in verbatim environments
\IfFileExists{upquote.sty}{\usepackage{upquote}}{}
\IfFileExists{microtype.sty}{% use microtype if available
  \usepackage[]{microtype}
  \UseMicrotypeSet[protrusion]{basicmath} % disable protrusion for tt fonts
}{}
\usepackage{setspace}
\makeatletter
\@ifundefined{KOMAClassName}{% if non-KOMA class
  \IfFileExists{parskip.sty}{%
    \usepackage{parskip}
  }{% else
    \setlength{\parindent}{0pt}
    \setlength{\parskip}{6pt plus 2pt minus 1pt}}
}{% if KOMA class
  \KOMAoptions{parskip=half}}
\makeatother
\ifLuaTeX
\usepackage[bidi=basic,shorthands=off]{babel}
\else
\usepackage[bidi=default,shorthands=off]{babel}
\fi
\ifPDFTeX
\else
\babelfont{rm}[]{Times New Roman}
\fi
\ifLuaTeX
  \usepackage{selnolig} % disable illegal ligatures
\fi
\setlength{\emergencystretch}{3em} % prevent overfull lines
\providecommand{\tightlist}{%
  \setlength{\itemsep}{0pt}\setlength{\parskip}{0pt}}
\usepackage{bookmark}
\IfFileExists{xurl.sty}{\usepackage{xurl}}{} % add URL line breaks if available
\urlstyle{same}
% fallback for those not using the hyperref driver hyperxmp:
\makeatletter
\@ifundefined{xmpquote}{\newcommand{\xmpquote}[1]{#1}}{}
\makeatother
\hypersetup{
  pdflang={pl-PL},
  colorlinks=true,
  linkcolor={Maroon},
  filecolor={Maroon},
  citecolor={Blue},
  urlcolor={Blue},
  pdfcreator={LaTeX via pandoc}}

\definecolor{heading}{HTML}{243447}
\definecolor{speech}{HTML}{F2F4F7}

\titleformat{\section}
  {\LARGE\bfseries\color{heading}}{}{0pt}{}
\titleformat{\subsection}
  {\Large\bfseries\color{heading}}{}{0pt}{}
\titleformat{\subsubsection}
  {\large\bfseries\color{heading}}{}{0pt}{}
\titlespacing*{\section}{0pt}{2.2em}{1em}
\titlespacing*{\subsection}{0pt}{1.8em}{0.8em}
\titlespacing*{\subsubsection}{0pt}{1.4em}{0.6em}

\renewenvironment{quote}
  {\list{}{\leftmargin=1.2em\rightmargin=0.6em}%
   \item\relax\color{black}\setlength{\parskip}{0.7em}}
  {\endlist}

\pagestyle{fancy}
\setlength{\headheight}{14pt}
\fancyhf{}
\fancyhead[L]{\small Scenariusz demonstracji aplikacji}
\fancyhead[R]{\small Obrona pracy inzynierskiej}
\fancyfoot[C]{\thepage}
\renewcommand{\headrulewidth}{0.4pt}

\author{}
\date{}

\begin{document}

\setstretch{1.2}
\section{Demonstracja aplikacji}\label{demonstracja-aplikacji}

\subsection{Plan czasowy}\label{plan-czasowy}

\begin{enumerate}
\def\labelenumi{\arabic{enumi}.}
\tightlist
\item
  Przeglad wszystkich funkcjonalnosci projektu - okolo 5 minut.
\item
  Dashboard, przypinanie, filtrowanie i sortowanie - okolo 4-5 minut.
\item
  Tworzenie i edycja notatki - okolo 3 minut.
\item
  Tworzenie i edycja checklisty - okolo 4 minut.
\item
  Podsumowanie - okolo 1 minuty.
\end{enumerate}

\section{Czesc I. Gotowa wypowiedz podczas
demonstracji}\label{czesc-i-gotowa-wypowiedz-podczas-demonstracji}

Podane czasy obejmuja wypowiedz i obsluge aplikacji. Wprowadzenie jest
czescia pieciominutowego przegladu funkcjonalnosci, dlatego nie stanowi
osobnej pozycji planu.

\subsection{1. Przeglad wszystkich funkcjonalnosci projektu - okolo 5
minut}\label{1-przeglad-wszystkich-funkcjonalnosci-projektu---okolo-5-minut}

\subsubsection{1.1. Krotkie wprowadzenie - okolo 40
sekund}\label{11-krotkie-wprowadzenie---okolo-40-sekund}

\textbf{Dzialanie:} Otworz dashboard. Sidebar powinien byc rozwiniety.

\textbf{Wypowiedz:}

\begin{quote}
Przedstawiona aplikacja sluzy do organizacji zadan oraz informacji
zwiazanych z ich realizacja. Uzytkownik moze tworzyc taski, checklisty,
notatki i projekty, a nastepnie laczyc je ze soba. Pozwala to przejsc od
prostej notatki do projektu zarzadzanego za pomoca tablicy Kanban bez
korzystania z kilku niezaleznych narzedzi.

Po lewej stronie znajduje sie wspolna nawigacja, a centralna czesc
widoku zmienia sie zaleznie od wybranego modulu. Dostep do danych jest
chroniony sesja uzytkownika utworzona po zalogowaniu przez zewnetrznego
dostawce OAuth.
\end{quote}

\subsubsection{1.2. Nawigacja i glowne typy elementow - okolo 50
sekund}\label{12-nawigacja-i-glowne-typy-elementow---okolo-50-sekund}

\textbf{Dzialanie:} Wskaz pozycje Dashboard, Tasks, Checklists, Notes,
Projects, Archive i Settings. Nie otwieraj jeszcze formularzy.

\textbf{Wypowiedz:}

\begin{quote}
Aplikacja zostala podzielona na glowne moduly widoczne w sidebarze.
Dashboard stanowi ekran startowy i agreguje najwazniejsze informacje.
Task jest podstawowym elementem reprezentujacym zadanie. Checklista
sklada sie z mniejszych czynnosci, ktore mozna kolejno oznaczac jako
wykonane. Notatki sluza do przechowywania swobodnej tresci i moga zostac
powiazane z innymi elementami. Projekty grupuja taski i checklisty oraz
udostepniaja widok Kanban. Osobny modul Archive przechowuje elementy
zarchiwizowane, a Settings odpowiada za wyglad aplikacji.

Na stronach list dostepne sa wspolne mechanizmy wyszukiwania,
filtrowania i sortowania. Elementy mozna takze ustawiac we wlasnej
kolejnosci przez drag and drop, a podobne operacje zachowuja sie
jednakowo w poszczegolnych modulach.
\end{quote}

\subsubsection{1.3. Taski i relacje - okolo 50
sekund}\label{13-taski-i-relacje---okolo-50-sekund}

\textbf{Dzialanie:} Otworz strone Tasks, pokaz kafelki i kontrolki
wyszukiwania. Otworz jeden przygotowany task, ale nie edytuj go
szczegolowo.

\textbf{Wypowiedz:}

\begin{quote}
Task zawiera tytul, opis, priorytet, termin, status, tagi oraz
opcjonalne powiazanie z projektem. Moze tez zawierac podlaczone
checklisty. Najwazniejsze pola, tytul i opis mozna edytowac bezposrednio
w widoku szczegolow.

Istotnym elementem jest zachowanie kontekstu nawigacji. Jezeli task
zostal otwarty z projektu albo z tablicy Kanban, przycisk Back oraz
przekierowanie po archiwizacji prowadza do poprzedniego widoku, a nie
zawsze do glownej listy taskow.
\end{quote}

\subsubsection{1.4. Projekty i Kanban - okolo 1 minuty i 50
sekund}\label{14-projekty-i-kanban---okolo-1-minuty-i-50-sekund}

\textbf{Dzialanie:} Otworz przygotowany projekt. Pokaz informacje
projektu, sekcje checklist i tablice Kanban. Przeciagnij jeden task w
obrebie kolumny, drugi pomiedzy kolumnami, a nastepnie przestaw dwie
kolumny.

\textbf{Wypowiedz:}

\begin{quote}
Projekt grupuje zadania dotyczace wspolnego celu. Taski mozna
przedstawic jako liste albo tablice Kanban. Kolumny maja wlasne nazwy,
kolory i kolejnosc; mozna je dodawac, edytowac, przestawiac oraz usuwac.

Przeniesienie taska do innej kolumny zmienia jego status, natomiast
przeniesienie w pionie aktualizuje pozycje w kolumnie. Linia pomocnicza
wskazuje miejsce wstawienia, a wynik operacji zostaje zapisany w bazie
danych.

Nad tablica znajduje sie sekcja checklist powiazanych z projektem.
Checkliste mozna rozwinac i obslugiwac bez opuszczania widoku projektu.
Mozna zmieniac stan jej elementow, edytowac ich nazwy, dodawac nowe
pozycje i zmieniac kolejnosc. Pozostaje ona jednoczesnie samodzielnym
dokumentem w bazie.
\end{quote}

\subsubsection{1.5. Archiwum i ustawienia - okolo 50
sekund}\label{15-archiwum-i-ustawienia---okolo-50-sekund}

\textbf{Dzialanie:} Krotko otworz Archive, a nastepnie Settings. Pokaz
akcje przywracania i trwalego usuwania bez ich zatwierdzania. W
ustawieniach zmien motyw i wroc do poprzedniego.

\textbf{Wypowiedz:}

\begin{quote}
Standardowe usuniecie elementu oznacza archiwizacje. Dokument znika z
aktywnych list, ale w archiwum mozna go przywrocic albo usunac trwale.
Obie operacje wymagaja potwierdzenia, co ogranicza ryzyko przypadkowej
utraty danych.

Aplikacja obsluguje motyw jasny, ciemny i systemowy. Dostepne sa motywy
domyslne oraz motywy utworzone przez uzytkownika. Kolory sa zapisywane
jako preferencje uzytkownika i przekazywane do interfejsu przez zmienne
CSS, dlatego zmiana jest widoczna w calej aplikacji.
\end{quote}

\subsection{2. Dashboard, przypinanie, filtrowanie i sortowanie - okolo
4-5
minut}\label{2-dashboard-przypinanie-filtrowanie-i-sortowanie---okolo-4-5-minut}

\textbf{Dzialanie:} Wroc do dashboardu.

\textbf{Wypowiedz:}

\begin{quote}
Dashboard jest ekranem zbiorczym aplikacji. Jego pierwsza czesc zawiera
skroty Quick create, dzieki ktorym mozna od razu przejsc do utworzenia
taska, checklisty, notatki albo projektu. Ponizej znajduja sie liczniki
poszczegolnych typow elementow oraz kalendarz oparty na terminach
zapisanych w taskach i projektach. Najechanie na licznik pozwala szybko
podejrzec ostatnio aktualizowane elementy.

Glowna czesc dashboardu jest przeznaczona na elementy przypiete przez
uzytkownika. Przypiecie nie kopiuje obiektu. W bazie powstaje osobny
dokument Pin zawierajacy typ i identyfikator elementu docelowego. Dzieki
temu na jednej tablicy moga wystepowac rozne typy danych, natomiast
oryginalne notatki, checklisty, taski i projekty nadal pozostaja w
swoich kolekcjach.
\end{quote}

\subsubsection{2.1. Przypinanie elementu}\label{21-przypinanie-elementu}

\textbf{Dzialanie:} Otworz przygotowana notatke albo checkliste. Kliknij
Pin. Wroc do dashboardu i wskaz nowy element. Nastepnie mozesz pokazac
Unpin.

\textbf{Wypowiedz:}

\begin{quote}
Kazdy element, ktory moze byc przypiety, ma przycisk Pin w widoku
szczegolow. Po przypieciu serwer zapisuje odniesienie do elementu, a
dashboard pobiera dokumenty pinow nalezace do aktualnego uzytkownika i
uzupelnia je danymi obiektow docelowych. Przycisk zmienia stan na Unpin,
dlatego ta sama kontrolka obsluguje obie operacje.

Piny mozna rowniez ustawiac metoda drag and drop. Pole position
przechowuje kolejnosc wybrana przez uzytkownika. Jest to domyslny sposob
prezentacji elementow, dopoki uzytkownik nie wybierze innego sortowania.
\end{quote}

\subsubsection{2.2. Wyszukiwanie i
filtrowanie}\label{22-wyszukiwanie-i-filtrowanie}

\textbf{Dzialanie:} Wpisz fragment tytulu przypietego elementu.
Nastepnie wyczysc tekst, wybierz typ, np. Checklists, i pokaz filtr
Relation: Linked, Linked to Project, Linked to Task oraz Not Linked.

\textbf{Wypowiedz:}

\begin{quote}
Wyszukiwanie dziala na biezaco, dlatego nie wymaga przycisku Apply.
Wpisanie fragmentu tytulu natychmiast ogranicza zestaw wynikow. Ponizej
znajduja sie filtry charakterystyczne dla dashboardu. Pierwszy pozwala
wybrac rodzaj elementu, a filtr Relation rozroznia elementy
niepowiazane, dowolnie powiazane, powiazane z projektem oraz powiazane z
taskiem.

Projekty nie sa ukrywane przez filtr relacji, poniewaz w modelu
aplikacji sa elementami nadrzednymi i same nie sa podpinane do innych
projektow lub taskow. Przycisk Reset przywraca wartosci domyslne
wszystkich kontrolek, lacznie z wyszukiwanym tekstem.
\end{quote}

\subsubsection{2.3. Sortowanie}\label{23-sortowanie}

\textbf{Dzialanie:} Pokaz opcje sortowania. Wybierz tytul albo date
aktualizacji, zmien Ascending/\allowbreak Descending, a na koncu wybierz
User\textquotesingle s Order.

\textbf{Wypowiedz:}

\begin{quote}
Elementy mozna sortowac wedlug dostepnych cech, na przyklad tytulu lub
daty aktualizacji. Dla takich kryteriow mozna wybrac kierunek rosnacy
albo malejacy. Opcja User\textquotesingle s Order korzysta natomiast z
pozycji zapisanej podczas drag and drop, dlatego kontrolka kierunku
sortowania znika - odwrocenie kolejnosci nie mialoby tu znaczenia
zgodnego z intencja uzytkownika.

Ten sam podstawowy komponent kontrolek jest wykorzystywany na stronach
notatek, checklist, taskow i projektow. Dashboard rozszerza go jedynie o
filtry typu elementu, poniewaz prezentuje dane z kilku kolekcji
jednoczesnie.
\end{quote}

\subsection{3. Tworzenie i edycja notatki - okolo 3
minut}\label{3-tworzenie-i-edycja-notatki---okolo-3-minut}

\subsubsection{3.1. Utworzenie notatki}\label{31-utworzenie-notatki}

\textbf{Dzialanie:} Uzyj Quick create lub przycisku \texttt{+\ New} na
stronie Notes. Wpisz przygotowany tytul, krotka tresc i dwa tagi. Utworz
notatke.

\textbf{Wypowiedz:}

\begin{quote}
Notatka jest najprostszym elementem aplikacji i sluzy do szybkiego
zapisania informacji. Formularz zawiera wymagany tytul, opcjonalna tresc
oraz tagi. Tagi mozna dynamicznie dodawac i usuwac. Przed wyslaniem dane
sa normalizowane, czyli usuwane sa puste tagi i zbedne spacje.

Po zatwierdzeniu formularz wysyla zadanie POST do endpointu notatek.
Backend sprawdza sesje, waliduje dane i zapisuje dokument z
identyfikatorem wlasciciela. Odpowiedz zawiera identyfikator nowej
notatki, dlatego aplikacja moze od razu przejsc do jej widoku szczegolow
zamiast wracac do ogolnej listy.
\end{quote}

\subsubsection{3.2. Edycja inline i
powiazania}\label{32-edycja-inline-i-powiazania}

\textbf{Dzialanie:} Kliknij tytul, zmien go i zakoncz edycje Enterem.
Kliknij tresc i dopisz zdanie. Edytuj tagi. Pokaz sekcje Linked items i
podlacz przygotowany task lub projekt. Zapisz zmiany.

\textbf{Wypowiedz:}

\begin{quote}
Widok szczegolow wykorzystuje edycje inline. Tekst w trybie podgladu
zmienia sie w pole dopiero po kliknieciu, przy czym zachowuje swoj
rozmiar, aby interfejs nie przesuwal sie podczas edycji. Enter
zatwierdza pojedyncza linie, a opuszczenie pola konczy edycje. Lokalny
stan komponentu przechowuje wersje robocza, natomiast przycisk Save
zapisuje komplet zmian przez zadanie PATCH.

Notatke mozna powiazac z innymi elementami aplikacji. Sekcja Linked
items zapisuje typ i identyfikator obiektu docelowego. Pozwala to
wykorzystac notatke jako dodatkowy kontekst dla taska, checklisty lub
projektu bez kopiowania jej tresci. Po zapisaniu relacja jest odtwarzana
przy kolejnym pobraniu dokumentu.

Z poziomu szczegolow notatke mozna przypiac do dashboardu albo przeniesc
do archiwum. Jezeli widok zostal otwarty z innego miejsca aplikacji,
parametr returnTo zachowuje informacje o poprzednim adresie i pozwala do
niego wrocic.
\end{quote}

\subsection{4. Tworzenie i edycja checklisty - okolo 4
minut}\label{4-tworzenie-i-edycja-checklisty---okolo-4-minut}

\subsubsection{4.1. Utworzenie
checklisty}\label{41-utworzenie-checklisty}

\textbf{Dzialanie:} Otworz formularz nowej checklisty. Wpisz tytul i co
najmniej trzy elementy. Zmien ich kolejnosc przez drag and drop, usun
jeden i dodaj go ponownie. Utworz checkliste.

\textbf{Wypowiedz:}

\begin{quote}
Checklista sklada sie z tytulu oraz tablicy elementow. Kazdy element
zawiera tresc, informacje o wykonaniu oraz pozycje. Formularz pozwala
dodawac i usuwac wiersze, a kolejnosc mozna zmieniac metoda drag and
drop jeszcze przed zapisaniem dokumentu.

Przed wyslaniem puste wiersze sa odrzucane, a pozostale elementy
otrzymuja kolejne wartosci position. Dane trafiaja do endpointu
checklist przez POST. Podobnie jak w przypadku notatki, po utworzeniu
aplikacja odczytuje identyfikator dokumentu i przechodzi bezposrednio do
jego szczegolow.
\end{quote}

\subsubsection{4.2. Realizacja i edycja
elementow}\label{42-realizacja-i-edycja-elementow}

\textbf{Dzialanie:} Zaznacz i odznacz element, klikajac caly wiersz.
Przenies element w inne miejsce. Uruchom edycje nazwy ikona olowka,
zmien tekst i zapisz. Dodaj nowy element.

\textbf{Wypowiedz:}

\begin{quote}
W widoku checklisty klikniecie w wiersz zmienia stan wykonania elementu.
Dla istniejacych pozycji ta operacja jest zapisywana od razu, poniewaz
uzytkownik oczekuje natychmiastowego zachowania statusu. Data wykonania
jest ustawiana przy zaznaczeniu i usuwana przy odznaczeniu.

Nazwa elementu jest edytowana dopiero po uzyciu ikony olowka. Zapobiega
to przypadkowemu uruchamianiu pola podczas zaznaczania checklisty. Drag
and drop nie zamienia tylko dwoch wartosci miejscami. Element jest
usuwany ze starego indeksu i wstawiany w nowe miejsce, a wszystkie
pozycje zostaja ponownie przeliczone.

Tytul i strukturalne zmiany checklisty korzystaja z wersji roboczej oraz
przycisku Save. Aplikacja porownuje stan poczatkowy z aktualnym, dzieki
czemu wie, czy wystepuja niezapisane zmiany. Uzytkownik moze tez uzyc
Cancel, aby przywrocic dane pobrane z serwera.
\end{quote}

\subsubsection{4.3. Checklista zagniezdzona w tasku lub
projekcie}\label{43-checklista-zagniezdzona-w-tasku-lub-projekcie}

\textbf{Dzialanie:} Jezeli pozostalo wystarczajaco czasu, otworz task
albo projekt z podlaczona checklista. Rozwin ja, zaznacz element, zmien
kolejnosc i wskaz ikony odlaczenia oraz otwarcia pelnego widoku.

\textbf{Wypowiedz:}

\begin{quote}
Ta sama checklista moze zostac podlaczona do taska albo projektu. W
widoku zagniezdzonym zachowuje podstawowe funkcje: rozwijanie,
zaznaczanie, edycje, dodawanie, usuwanie i zmiane kolejnosci. Zmiany sa
zapisywane bez osobnego przycisku Save.

Komponent ogranicza jednak liczbe zbednych zapytan. Wpisywanie nazwy nie
wysyla zadania po kazdym znaku - zapis odbywa sie po zakonczeniu edycji.
Dodatkowo zapamietywany jest ostatni zapisany stan checklisty. Jezeli
dane sie nie zmienily, kolejne zadanie PATCH nie jest wykonywane.
\end{quote}

\subsection{5. Podsumowanie - okolo 1
minuty}\label{5-podsumowanie---okolo-1-minuty}

\textbf{Dzialanie:} Wroc do dashboardu, aby zakonczyc prezentacje na
ekranie zbiorczym.

\textbf{Wypowiedz:}

\begin{quote}
Podsumowujac, aplikacja laczy proste elementy organizacyjne z bardziej
rozbudowanym zarzadzaniem projektami. Uzytkownik moze rozpoczac od
notatki lub checklisty, utworzyc task, polaczyc te dane z projektem, a
nastepnie kontrolowac realizacje na tablicy Kanban. Dashboard zapewnia
szybki dostep do najwazniejszych informacji, a wspolne filtrowanie,
sortowanie, przypinanie i archiwizacja utrzymuja spojne zachowanie
calego systemu.

Od strony technicznej aplikacja laczy frontend React i Next.js z
backendowymi Route Handlers oraz baza MongoDB obslugiwana przez
Mongoose. Operacje sa chronione sesja uzytkownika, dane sa walidowane
przed zapisem, a kluczowe formularze, endpointy i mechanizmy interfejsu
zostaly objete testami automatycznymi.
\end{quote}

\end{document}
